\documentclass[11pt,a4paper]{article}
\usepackage[T1]{fontenc}
\usepackage[utf8]{inputenc}
\usepackage[provide=*,lithuanian]{babel}
\usepackage{lmodern,amsmath,amssymb,enumitem}
\usepackage[margin=1.5cm]{geometry}
\setlength{\parindent}{0pt}
\setlength{\emergencystretch}{3em}
\setlength{\abovedisplayskip}{4pt}
\setlength{\belowdisplayskip}{4pt}
\setlist[enumerate]{label=\textbf{\arabic*.},leftmargin=0.7cm,itemsep=0.05cm,topsep=0.2cm,parsep=0pt}
\begin{document}
\begin{center}
{\Large\bfseries Vijetos teoremos taikymas}

\medskip
{\large Lygis: Beast}

\end{center}
\textbf{Užduotis.} Atlikite uždavinius, naudodamiesi Vijetos teorema.
Uždaviniuose su parametru patikrinkite, ar gautos jo reikšmės užtikrina
realiųjų šaknų egzistavimą.
\medskip
\begin{enumerate}
\item Duotas daugianaris \(P(x)=x^2-4x+1\), kurio šaknys yra \(x_1,x_2\). Nerasdami šaknų, apskaičiuokite \mbox{\(x_1^3+x_2^3\)}.
\item Lygties \(x^2-5x+3=0\) šaknys yra \(x_1,x_2\). Nerasdami šaknų, apskaičiuokite \mbox{\(\frac1{x_1-1}+\frac1{x_2-1}\)}.
\item Lygties \(x^2-3x-2=0\) šaknys yra \(x_1,x_2\). Nerasdami šaknų, apskaičiuokite \mbox{\(x_1^3x_2+x_1x_2^3\)}.
\item Raskite visas \(m\) reikšmes, su kuriomis daugianario \(P(x)=x^2-(m+2)x+2m\) realiosios šaknys \(x_1,x_2\) tenkina \mbox{\(x_1^2+x_2^2=20\)}.
\item \begin{minipage}[t]{\linewidth}Kvadratinio daugianario \(P(x)\) vyriausiasis koeficientas lygus \(1\). Jo realiosios šaknys tenkina sąlygas: \[x_1^2+x_2^2=29,\quad x_1x_2=10,\quad x_1+x_2>0.\] Raskite \(P(x)\).\end{minipage}
\item Duoto daugianario \(P(x)=x^2-3x-2\) šaknys yra \(x_1,x_2\). Nerasdami jų, sudarykite kvadratinę lygtį, kurios koeficientas prie \(x^2\) lygus \(1\), o šaknys yra \mbox{\(x_1+2,\quad x_2+2\)}.
\item Lygties \(x^2-3x+1=0\) šaknys yra \(x_1,x_2\). Nerasdami šaknų, apskaičiuokite \mbox{\(x_1^4+x_2^4\)}.
\item Daugianario \(P(x)=x^2-6x+m\) viena realioji šaknis lygi kitos šaknies kvadratui. Raskite visas galimas \(m\) reikšmes.
\item Duotas daugianaris \(P(x)=x^2-5x+2\), kurio šaknys yra \(x_1,x_2\). Nerasdami šaknų, apskaičiuokite \mbox{\(\frac{x_1}{x_2}+\frac{x_2}{x_1}\)}.
\item Raskite visas \(m\) reikšmes, su kuriomis daugianario \(P(x)=x^2-(m+1)x+2m\) realiosios šaknys skiriasi \(1\).
\item Lygties \(x^2+3x-2=0\) šaknys yra \(x_1,x_2\). Nerasdami šaknų, apskaičiuokite \mbox{\(x_1^3+x_2^3\)}.
\item Lygties \(x^2-3x+1=0\) šaknys yra \(x_1,x_2\). Nerasdami šaknų, apskaičiuokite \mbox{\(\frac{x_1^2}{x_2}+\frac{x_2^2}{x_1}\)}.
\item Duotas daugianaris \(P(x)=x^2-2x-1\), kurio šaknys yra \(x_1,x_2\). Nerasdami šaknų, apskaičiuokite \mbox{\((x_1^2+x_2^2)^2\)}.
\item Duotas daugianaris \(P(x)=x^2+2x-3\), kurio šaknys yra \(x_1,x_2\). Nerasdami šaknų, apskaičiuokite \mbox{\(x_1^2x_2+x_1x_2^2\)}.
\item Duotas daugianaris \(P(x)=x^2+3x-1\), kurio šaknys yra \(x_1,x_2\). Nerasdami šaknų, apskaičiuokite \mbox{\(\frac1{x_1+2}+\frac1{x_2+2}\)}.
\item Lygties \(x^2+4x+1=0\) šaknys yra \(x_1,x_2\). Nerasdami jų, sudarykite kvadratinę lygtį, kurios koeficientas prie \(x^2\) lygus \(1\), o šaknys yra \mbox{\(x_1-1,\quad x_2-1\)}.
\item \begin{minipage}[t]{\linewidth}Sudarykite kvadratinę lygtį, kurios koeficientas prie \(x^2\) lygus \(1\), o realiosios šaknys tenkina: \[x_1^2+x_2^2=13,\quad x_1x_2=-2,\quad x_1+x_2<0.\]\end{minipage}
\item Lygties \(x^2-mx+2m-2=0\) realiosios šaknys yra \(x_1,x_2\). Raskite visas \(m\) reikšmes, su kuriomis \mbox{\(x_1^2+x_2^2=16\)}.
\item Raskite visas \(m\) reikšmes, su kuriomis lygties \(x^2-2x+m-1=0\) realiosios šaknys gali būti pažymėtos taip, kad \mbox{\(x_1=x_2^2\)}.
\item Raskite visas \(m\) reikšmes, su kuriomis lygties \(x^2-(m+2)x+m+1=0\) šaknims \(x_1<x_2\) galioja \mbox{\(x_2-x_1=2\)}.
\item Duotas daugianaris \(P(x)=x^2-x-1\), kurio šaknys yra \(x_1,x_2\). Nerasdami šaknų, apskaičiuokite \mbox{\(\frac{x_1^2}{x_2^2}+\frac{x_2^2}{x_1^2}\)}.
\item Daugianario \(P(x)=4x^2-3x+2m\) viena realioji šaknis lygi kitos šaknies kvadratui. Raskite visas galimas \(m\) reikšmes.
\item Lygties \(2x^2-3x-1=0\) šaknys yra \(x_1,x_2\). Nerasdami šaknų, apskaičiuokite \mbox{\(\frac1{2x_1-1}+\frac1{2x_2-1}\)}.
\item Raskite visas \(m\) reikšmes, su kuriomis daugianario \(P(x)=x^2-2mx+m+1\) realiosios šaknys skiriasi \(2\).
\item \begin{minipage}[t]{\linewidth}Kvadratinio daugianario \(P(x)\) vyriausiasis koeficientas lygus \(1\). Jo realiosios šaknys tenkina sąlygas: \[(x_1-x_2)^2=9,\quad x_1x_2=4,\quad x_1+x_2>0.\] Raskite \(P(x)\).\end{minipage}
\item Lygties \(x^2-4x+2=0\) šaknys yra \(x_1,x_2\). Nerasdami šaknų, apskaičiuokite \mbox{\(x_1^3x_2^2+x_1^2x_2^3\)}.
\item Lygties \(x^2-4x+2=0\) šaknys yra \(x_1,x_2\). Nerasdami šaknų, apskaičiuokite \mbox{\((x_1^2-x_2^2)^2\)}.
\item Duotas daugianaris \(P(x)=x^2-5x+2\), kurio šaknys yra \(x_1,x_2\). Nerasdami šaknų, apskaičiuokite \mbox{\((x_1+x_2)^3-x_1^3-x_2^3\)}.
\item Duoto daugianario \(P(x)=2x^2-6x+1\) šaknys yra \(x_1,x_2\). Nerasdami jų, sudarykite kvadratinę lygtį, kurios koeficientas prie \(x^2\) lygus \(1\), o šaknys yra \mbox{\(2x_1,\quad 2x_2\)}.
\item Raskite visas \(m\) reikšmes, su kuriomis daugianario \(P(x)=2x^2-(m+4)x+2m\) realiosios šaknys \(x_1,x_2\) tenkina \mbox{\(x_1^2+x_2^2=13\)}.
\item Raskite visas \(m\) reikšmes, su kuriomis lygties \(2x^2-(m+2)x+m=0\) šaknims \(x_1<x_2\) galioja \mbox{\(x_2-x_1=\frac{1}{2}\)}.
\item Lygties \(x^2+2x-3=0\) šaknys yra \(x_1,x_2\). Nerasdami šaknų, apskaičiuokite \mbox{\((x_1+x_2)\left(\frac{x_1}{x_2}+\frac{x_2}{x_1}\right)\)}.
\item \begin{minipage}[t]{\linewidth}Sudarykite kvadratinę lygtį, kurios koeficientas prie \(x^2\) lygus \(1\), o realiosios šaknys tenkina: \[x_1+x_2=-5,\quad x_1^2+x_2^2=17.\]\end{minipage}
\item Duotas daugianaris \(P(x)=2x^2-6x+1\), kurio šaknys yra \(x_1,x_2\). Nerasdami šaknų, apskaičiuokite \mbox{\(x_1^3x_2+x_1x_2^3+2x_1^2x_2^2\)}.
\item Duotas daugianaris \(P(x)=x^2-x-2\), kurio šaknys yra \(x_1,x_2\). Nerasdami šaknų, apskaičiuokite \mbox{\(x_1^4+x_2^4-2x_1^2x_2^2\)}.
\item Raskite visas \(m\) reikšmes, su kuriomis lygties \(x^2-mx-8=0\) realiosios šaknys gali būti pažymėtos taip, kad \mbox{\(x_1=x_2^2\)}.
\item Lygties \(-x^2+2x+3=0\) šaknys yra \(x_1,x_2\). Nerasdami šaknų, apskaičiuokite \mbox{\(x_1^3+x_2^3+3x_1x_2(x_1+x_2)\)}.
\item Lygties \(-x^2+5x+2=0\) šaknys yra \(x_1,x_2\). Nerasdami jų, sudarykite kvadratinę lygtį, kurios koeficientas prie \(x^2\) lygus \(1\), o šaknys yra \mbox{\(-x_1,\quad -x_2\)}.
\item Duotas daugianaris \(P(x)=x^2-4x+2\), kurio šaknys yra \(x_1,x_2\). Nerasdami šaknų, apskaičiuokite \mbox{\(\frac{x_1}{x_1+1}+\frac{x_2}{x_2+1}\)}.
\item Lygties \(-x^2+(m-1)x+m+1=0\) realiosios šaknys yra \(x_1,x_2\). Raskite visas \(m\) reikšmes, su kuriomis \mbox{\(x_1^2+x_2^2=12\)}.
\item Duoto daugianario \(P(x)=x^2-5x+3\) šaknys yra \(x_1,x_2\). Nerasdami jų, sudarykite kvadratinę lygtį, kurios koeficientas prie \(x^2\) lygus \(3\), o šaknys yra \mbox{\(\frac1{x_1},\quad\frac1{x_2}\)}.
\item Duotas daugianaris \(P(x)=x^2-3x+1\), kurio šaknys yra \(x_1,x_2\). Nerasdami šaknų, apskaičiuokite \mbox{\((x_1+1)^3+(x_2+1)^3\)}.
\item Duotas daugianaris \(P(x)=2x^2-6x+4\), kurio šaknys yra \(x_1,x_2\). Nerasdami šaknų, apskaičiuokite \mbox{\(\left(\frac{x_1}{x_2}-\frac{x_2}{x_1}\right)^2\)}.
\item Lygties \(-x^2+3x+1=0\) šaknys yra \(x_1,x_2\). Nerasdami šaknų, apskaičiuokite \mbox{\(x_1^4x_2+x_1x_2^4\)}.
\item Daugianario \(P(x)=2x^2-(m+1)x+16\) viena realioji šaknis lygi kitos šaknies kvadratui. Raskite visas galimas \(m\) reikšmes.
\item Lygties \(x^2+2x-2=0\) šaknys yra \(x_1,x_2\). Nerasdami šaknų, apskaičiuokite \mbox{\(\frac{x_1+1}{x_1-1}+\frac{x_2+1}{x_2-1}\)}.
\item \begin{minipage}[t]{\linewidth}Kvadratinio daugianario \(P(x)\) vyriausiasis koeficientas lygus \(2\). Jo realiosios šaknys tenkina sąlygas: \[x_1^2+x_2^2=2,\quad x_1x_2=-\frac12,\quad x_1+x_2<0.\] Raskite \(P(x)\).\end{minipage}
\item Lygties \(x^2+2x-1=0\) šaknys yra \(x_1,x_2\). Nerasdami šaknų, apskaičiuokite \mbox{\(x_1^4+x_2^4+x_1^2x_2^2\)}.
\item Raskite visas \(m\) reikšmes, su kuriomis daugianario \(P(x)=x^2-(2m+1)x+m\) realiosios šaknys \(x_1,x_2\) tenkina \mbox{\(x_1^2+x_2^2=7\)}.
\item Raskite visas \(m\) reikšmes, su kuriomis daugianario \(P(x)=-x^2+(m-1)x-2m+6\) realiosios šaknys skiriasi \(3\).
\end{enumerate}
\clearpage
\section*{Atsakymai}
\begingroup\small
\setlength{\baselineskip}{11pt}
\setlist[enumerate]{itemsep=0cm}
\begin{enumerate}
\item \(52\).
\item \(-3\).
\item \(-26\).
\item \(m_1=-4,\quad m_2=4\).
\item \(P(x)=x^2-7x+10\).
\item \(x^2-7x+8=0\).
\item \(47\).
\item \(m_1=-27,\quad m_2=8\).
\item \(\frac{21}{2}\).
\item \(m_1=0,\quad m_2=6\).
\item \(-45\).
\item \(18\).
\item \(36\).
\item \(6\).
\item \(-\frac{1}{3}\).
\item \(x^2+6x+6=0\).
\item \(x^2+3x-2=0\).
\item \(m_1=-2\).
\item \(m_1=-7,\quad m_2=2\).
\item \(m_1=-2,\quad m_2=2\).
\item \(7\).
\item \(m_1=-\frac{27}{4},\quad m_2=\frac{1}{4}\).
\item \(-\frac{1}{4}\).
\item \(m_1=-1,\quad m_2=2\).
\item \(P(x)=x^2-5x+4\).
\item \(16\).
\item \(128\).
\item \(30\).
\item \(x^2-6x+2=0\).
\item \(m_1=-6,\quad m_2=6\).
\item \(m_1=1,\quad m_2=3\).
\item \(\frac{20}{3}\).
\item \(x^2+5x+4=0\).
\item \(\frac{9}{2}\).
\item \(9\).
\item \(m_1=2\).
\item \(8\).
\item \(x^2+5x-2=0\).
\item \(\frac{8}{7}\).
\item \(m_1=-3,\quad m_2=3\).
\item \(3x^2-5x+1=0\).
\item \(50\).
\item \(\frac{9}{4}\).
\item \(-36\).
\item \(m_1=11\).
\item \(-6\).
\item \(P(x)=2x^2+2x-1\).
\item \(35\).
\item \(m_1=-\frac{3}{2},\quad m_2=1\).
\item \(m_1=2,\quad m_2=8\).
\end{enumerate}
\endgroup
\end{document}
